\subsection{单射道路}

命题 18. --- 设 X 是一个豪斯多夫拓扑空间。设 a 和 b 是 X 中属于同一道路连通分支的彼此不同的点。存在一条在 X 中连接 a 与 b 的单射道路。

设 \(f \colon \mathbf{I} \to \mathbf{X}\) 是连接 \(a\) 与 \(b\) 的 \(\mathbf{X}\) 中的一条道路。令 \(\mathcal{U}\) 为由开子集 \(\mathrm{U}\) 组成的集合，这些子集属于 {]}0,1{[}，并满足 \(f(x) = f(y)\) 对 U 的每个连通分支 {]}x, y{[} 成立。

引理 8. --- 按包含关系排序的集合 U 是归纳的。

设 V 是 U 的一个全序子集。我们来证明属于 V 的各集合的并 V 是 U 的一个元素，即，对于 V 的每个连通分支 {]}u, v{[}，都有 \(f(u) = f(v)\)。

令 \(x\) 为 \([u, v[\) 中的一点。令 \(\mathcal{V}_x\) 为这样的 \(U \in \mathcal{V}\) 组成的集合：\(x \in U\)；对于这样的 \(U\)，记 \([u_U, v_U[\) 为 \(x\) 在 \(U\) 中的连通分支。有 \(u_{U'} \leqslant u_U < v_U \leqslant v_{U'}\)，若 \(U\) 和 \(U'\) 是 \(\mathcal{V}_x\) 中满足 \(U \subset U'\) 的元素。由 III, p.~280 的引理 7，所有 \(U \in \mathcal{V}_x\) 的 \([u_U, v_U[\) 的并等于 {]}u, v{[}，故有 u = lim \(u_{U}\) 和 v = lim \(v_{U}\)，其中极限取遍有向集 \(V_{x}\)。由于映射 f 连续且拓扑空间 X 是豪斯多夫的，等式 \(f(u_{\mathrm{U}}) = f(v_{\mathrm{U}})\) 对所有 \(U \in V_{x}\) 蕴含 \(f(u) = f(v)\)，这正是所要证明的。

根据 E, III, p.~20, 定理 2，存在属于 \(\mathcal{U}\) 的一个开子集 U，它关于包含关系是极大的。

令 \(g \colon \mathbf{I} \to \mathrm{X}\) 按如下方式定义。若 \(t \notin \mathrm{U}\)，则置 \(g(t) = f(t)\)；否则，记 \(]u, v[\) 为 \(t\) 在 \(\mathrm{U}\) 中的连通分支，并置 \(g(t) = f(u) = f(v)\)，于是映射 \(g|[u, v]\) 是像为 \(f(u)\) 的常值映射。

命题 18 由下述引理推出。

引理 9. --- 存在连续的递增满射 \(u\colon\mathbf{I}\to\mathbf{I}\) 以及单射道路 \(c\colon\mathbf{I}\to\mathbf{X}\)，它连接 \(a\) 与 \(b\)，并使得 \(g=c\circ u\)。

我们先证明映射 \(g\) 连续。令 \(t\) 为 \(\mathbf{I}\) 中的一点。若 \(t \in \mathrm{U}\)，则 \(g\) 在 \(t\) 的一个邻域内为常值，因而在 \(t\) 处连续。假设 \(t \notin \mathrm{U}\)，令 \(\mathrm{W}\) 为 \(g(t)\) 在 \(\mathrm{X}\) 中的一个开邻域。令 \(\mathrm{V}\) 为 \(\mathbf{I}\) 中包含 \(t\) 且满足 \(f(\mathrm{V}) \subset \mathrm{W}\) 的一个开区间；这样的区间存在，因为 \(f\) 在 \(t\) 处连续。我们来证明 \(g(\mathrm{V}) \subset \mathrm{W}\)。令 \(x \in \mathrm{V}\)；若 \(x \notin \mathrm{U}\)，则有 \(g(x) = f(x)\)，因而 \(g(x) \in \mathrm{W}\)。于是设 \(x \in \mathrm{U}\)，并令 \(]u, v[\) 为 \(x\) 在 \(\mathrm{U}\) 中的连通分支。注意 \(]u, v[\) 不包含 \(t\)。因此，若 \(x < t\)，则有 \(x < v \leqslant t\)，从而 \(v \in \mathrm{V}\) 且 \(g(x) = g(v) = f(v) \in f(\mathrm{V}) \subset \mathrm{W}\)。同样可证 \(g(x) \in \mathrm{W}\)，若 \(x > t\)。于是 \(g\) 在 \(t\) 处连续。

设 \(u\) 和 \(v\) 是 \(\mathbf{I}\) 中满足 \(g(u) = g(v)\) 且 \(u < v\) 的点。我们来证明 \(]u, v[ \subset U\)。置 \(U' = U \cup ]u, v[\)；它是 \(]0, 1[\) 的一个开子集。

若 \(u\) 属于 U，记 \(u'\) 为 \(u\) 在 U 中的连通分支在 I 中的最大下界；置 \(u' = u\)，若 \(u\) 不属于 U。同样，若 \(v\) 属于 U，记 \(v'\) 为 \(v\) 在 U 中的连通分支在 I 中的最小上界；置 \(v' = v\)，若 \(v\) 不属于 U。于是 \(]u', v'[\) 是 U' 的一个连通分支，且有 \(f(u') = g(u) = g(v) = f(v')\)。由于 U' 中异于 \(]u', v'[\) 的连通分支都是 U 的连通分支，开集 U' 是 \(\mathcal{U}\) 的一个元素。由于 U 关于包含关系是 \(\mathcal{U}\) 的极大元素，有 U' = U，且 \(]u, v[\) 包含于 U。这证明了映射 \(g\) 在区间 \([u, v]\) 上为常值。因此，\(g\) 的纤维是 I 的区间；由于 g 连续，这些区间是闭的。

记 R 为与 g 相伴的等价关系，记 p 为从 I 到 I/R 的典范满射。存在唯一的连续映射 \(g'\) 从 I/R 到 X，使得 \(g = g' \circ p\)；该映射是单射。由 III, p.~276 的命题 16 的推论，存在递增、连续且满射的映射 u，使得 R 是与 u 相伴的等价关系。由于空间 I 是紧的而空间 X 是豪斯多夫的，映射 u 是闭映射，因而是严格的（I, p.~18, 例 2）。它通过取商定义了从 I/R 到 I 的同胚 \(u'\)。于是，从 I 到 X 的映射 \(g' \circ (u')^{-1}\) 是一条起点为 a、终点为 b 的单射道路。

